\documentclass[12pt]{article}
\usepackage[utf8]{inputenc}
\usepackage[T1]{fontenc}
\usepackage[margin=1in]{geometry}
\usepackage{setspace}
\usepackage{xcolor}
\usepackage[colorlinks=true, urlcolor=blue!60!black, linkcolor=blue!60!black]{hyperref}
\setlength{\parindent}{0em}
\setlength{\parskip}{0.6em}
\renewcommand{\baselinestretch}{1.0}
\begin{document}
\begin{center}
{\LARGE \textit{Home Building \& Loan Ass'n v. Blaisdell}, 290 U.S. 398 (1934)\par}
\vspace{0.5em}
{\large Case Brief --- Two-Page Edition\par}
\end{center}

\section{Memory Jogger}
The Contracts Clause is not absolute; a State may make a reasonable, temporary, conditioned use of its police power in an emergency, as Minnesota did with its Depression-era mortgage moratorium (\href{https://supreme.justia.com/cases/federal/us/290/398/}{Justia}).

\section{Detailed Case Facts}
Minnesota's 1933 Mortgage Moratorium Law let courts extend redemption periods after foreclosure, no later than May 1, 1935, if the mortgagor paid the property's reasonable rental value toward taxes, insurance, and interest. The Blaisdells' home had been foreclosed and bought by the lender.

\section{Procedural History}
After an initial reversal, the trial court set rent at \$40 a month and extended redemption to May 1, 1935. The Minnesota Supreme Court affirmed, and the lender appealed.

\section{Judicial Votes}
Affirmed, 5--4. Hughes for the Court (Brandeis, Stone, Roberts, Cardozo); Sutherland, with Van Devanter, McReynolds, and Butler, dissenting.

\section{Holding}
The Contracts Clause must be harmonized with the State's reserved police power. Temporary, emergency relief serving a public end on reasonable conditions that protect the creditor's investment is valid; the law also satisfied due process and equal protection.

\section{Analysis of Opinions}
\textbf{Hughes:} An emergency does not create power but may furnish the occasion to use it. The question is whether the law serves a legitimate end by reasonable and appropriate means. A genuine emergency, public purpose, tailored judicial relief, preservation of the debt and rent, and temporary operation saved the law; the Clause's ``great clauses'' are read in light of the Nation's ``whole experience.'' \textbf{Sutherland:} the Clause was written precisely to forbid debtor relief in depressions, and its meaning does not change with economic conditions.

\section{Examples: Future Applications}
\textbf{Same-Side:} A temporary, rent-conditioned pandemic foreclosure pause is valid. \textbf{Opposite-Side:} Permanent, uncompensated commercial rent forgiveness fails, as in \textit{W. B. Worthen Co. v. Kavanaugh}. \textbf{Fence-Sitter:} Repealing a retirees' cost-of-living adjustment turns on whether a contract existed and on \textit{U.S. Trust} scrutiny.

\section{Critique}
\textit{Blaisdell} pits living constitutionalism against originalism on a clause aimed at stay laws, and shows how emergency exceptions harden into routine deference.

\section{Key Quotations}
\begin{quote}
The Contracts Clause's prohibition ``is not an absolute one, and is not to be read with literal exactness, like a mathematical formula.''
\end{quote}
\noindent --- Chief Justice Hughes, 290 U.S. at 428.

\end{document}
